\begin{proposition}
	Пусть $\Gamma$ "--- обобщённая окружность, точки $z_1, z_2 \in \CM$ симметричны относительно $\Gamma$, и $w(z)$ "--- дробно-линейное отображение. Тогда точки $w_1 := w(z_1)$ и $w_2  := w(z_2)$ симметричны относительно $\widetilde\Gamma := w(\Gamma)$.
\end{proposition}

\begin{proof}
	Пусть сначала $z_1 \ne z_2$. Рассмотрим произвольную обобщённую окружность $\widetilde{\gamma}$, проходящую через $w_1$ и $w_2$. Тогда $\gamma := w^{-1}(\widetilde\gamma)$ "--- обобщённая окружность, проходящая через $z_1$ и $z_2$ и потому пересекающая $\Gamma$ под прямым углом. Но отображение $w(z)$ конформно на $\CM$, поэтому $\widetilde \gamma$ тоже пересекает $\widetilde \Gamma$ под прямым углом. В силу произвольности выбора обобщённой окружности $\widetilde\gamma$, получено требуемое. Если же $z_1 = z_2$, то $z_1 = z_2 \in \Gamma$, откуда $w_1 = w_2 \in \widetilde\Gamma$, и симметричность этих точек очевидна.
\end{proof}

\begin{proposition}
	Пусть $D := \{z \in \Cm: |z| < 1\}$, $G := \{w \in \Cm: |w| < 1\}$, и функция $f$ конформно отображает $D$ на $G$. Тогда существуют $z_0 \in D$ и $\alpha \in [0, 2\pi)$ такие, что $f$ имеет следующий вид:
	\[f(z) = \frac{z - z_0}{1 - \overline{z_0}z}e^{i\alpha},~z \in D\]
\end{proposition}

\begin{proof}
	Зафиксируем $z_0 := f^{-1} \in D$ и рассмотрим следующее отображение:
	\[w(z) := \frac{z - z_0}{1 - z\overline{z_0}}\]
	
	Параметризуем окружность $\partial D$ как $z = e^{it}$, $t \in [0, 2\pi)$. Тогда для любого $t \in [0, 2\pi)$ выполнены следующие равенства:
	\[\left|w\left(e^{it}\right)\right| = \left|\frac{e^{it} - z_0}{1 - e^{it}\overline{z_0}}\right| = \frac{|e^{it} - z_0|}{|e^{-it} - \overline{z_0}|} = \frac{|e^{it} - z_0|}{|\overline{e^{it} - z_0}|} = 1\]
	
	Значит, $w(\partial D) \subset \partial G$, причем имеет место равенство  $w(\partial D) = \partial G$ в силу кругового свойства. Проверим, что тогда $w(D) \subset G$. Действительно, пусть существует такая точка $z_1 \in D$, что $w(z_1) \not\in G$, тогда отрезок $[z_0, z_1]$ переходит в дугу с концами $w(z_0) = 0$ и $w(z_1)$, поэтому существует точка $z^* \in [z_0, z_1]$ такая, что $w(z^*) \in \partial G$. Но это невозможно, поскольку $z^* \not \in \partial D$, а отображение $w(z)$ конформно. Аналогично доказывается, что $w^{-1}(G) \subset D$, тогда $w(D) = G$.
	
	Итак, $w(z)$ конформно отображает $D$ на $G$, тогда и отображения вида $g(z) = w(z)e^{i\alpha}$ для произвольного $\alpha \in [0, 2\pi)$. В частности, это верно для $\alpha := \arg{f'(z_0)}$. Но, по теореме Римана, отображение $f$ однозначно задается значениями $z_0 = f^{-1}(0)$ и $\alpha = \arg{f'(z_0)}$, поэтому оно имеет требуемый вид.
\end{proof}

\begin{proposition}
	Пусть $D := \{z \in \Cm: \im{z} > 0\}$, $G := \{w \in \Cm: |w| < 1\}$, и функция $f$ конформно отображает $D$ на $G$. Тогда существуют $z_0 \in D$ и $\alpha \in [0, 2\pi)$ такие, что $f$ имеет следующий вид:
	\[f(z) = \frac{z - z_0}{z - \overline{z_0}}e^{i\alpha},~z \in D\]
\end{proposition}

\begin{proof}
	Зафиксируем $z_0 \in D$ и рассмотрим следующее отображение:
	\[w(z) := \frac{z - z_0}{z - \overline{z_0}}\]
	
	Параметризуем прямую $\partial D$ как $z = t$, $t \in \R$. Тогда для любого $t \in \R$ выполнены следующие равенства:
	\[\left|w\left(t\right)\right| = \left|\frac{t - z_0}{t - \overline{z_0}}\right| = \frac{|t - z_0|}{\left|\overline{t - z_0}\right|} =  1\]
	
	Значит, $w(\partial D) \subset \partial G$, причем имеет место равенство  $w(\partial D) = \partial G$ в силу кругового свойства. Дальнейшее рассуждение аналогично предыдущему утверждению, но в данном случае следует положить $\alpha := \arg{f'(z_0)} + \frac\pi 2$.
\end{proof}